% !Mode:: "TeX:UTF-8"
\chapter*{结论\markboth{结论}{}}
\addcontentsline{toc}{chapter}{结论}

本文针对城市交通管理中自然语言交互与时空数据智能分析的实际需求，设计并实现了一个基于大语言模型的交通智能体系统。该系统以 Qwen3.5-9B 为核心推理引擎，以 PDFormer 为底层时空预测模型，以 LangGraph 为状态图调度框架，通过慢思考与快反应相结合的双层架构，打通了从自然语言指令到专业交通分析报告的完整链路。本文的主要工作总结如下。

第一，设计了一种面向交通领域的智能体双层驱动架构。该架构由慢思考层和快反应层组成。慢思考层由 Qwen3.5-9B 大模型负责理解用户意图、选择工具并生成调用指令。快反应层由工具执行节点负责并行调用底层数据服务。两层之间通过 LangGraph 的 ReAct 推理循环进行编排，实现了推理与执行的有效分离。在工具调用的解析上，系统采用了 LangChain 原生 bind\_tools 优先、正则表达式回退的双重兼容机制，兼顾了标准 API 通道的可靠性和原生 XML 格式的灵活性。

第二，完成了大模型的本地化部署与工具链集成。本研究将 Qwen3.5-9B 原始精度模型通过 vLLM 推理框架部署在本地 GPU 服务器上，支持 20K Token 的上下文窗口。围绕 PeMS08 交通数据集，设计并实现了三个核心工具：simple\_query 负责单点交通状态查询，deep\_analysis 负责结合历史趋势进行深度研判，predict\_traffic 负责调用 PDFormer 模型对未来流量进行外推预测。三个工具通过 FastAPI 微服务架构封装，与智能体决策层之间通过标准 HTTP 协议进行解耦通信。

第三，设计并执行了一套覆盖七类场景的智能体测试方案。基于 37 条测试用例，从意图识别、工具调度、报告质量和异常处理四个维度对智能体进行了系统性评估。实验结果表明，智能体在所有 37 条测试用例中均正确识别了用户意图，能够适应五种不同的自然语言时间格式，在传感器查询、拥堵研判、流量预测和多工具协同等场景下均能生成结构化的专业分析报告。在异常处理方面，智能体对传感器不存在、时间超范围和历史数据不足等异常情况均做出了有意义的降级处理，未出现崩溃或幻觉。

第四，在实验中发现了多工具协同场景下的数字歧义问题。当传感器编号落在 1–31 的合法日期范围内且紧邻时间描述词出现时，大模型可能将传感器编号误解为日期，从而干扰时间参照的解析。传感器编号大于 31 时该问题不再出现。这一发现对后续的提示词工程优化具有参考价值。

本研究的不足之处主要有三个方面。第一，PeMS08 数据集仅覆盖 2016 年 7–8 月的历史数据，限制了智能体对相对时间和实时查询的支持能力。第二，当前仅实现了三个核心工具，尚未覆盖路径规划、信号灯配时优化等更丰富的交通管理功能。第三，受限于单卡 GPU 的显存容量，模型参数规模和推理吞吐量仍有提升空间。

在未来的工作中，可以从以下几个方面继续推进。第一，接入实时交通数据流，将系统从历史数据分析升级为真正的在线交通感知与决策系统。第二，扩展工具链，引入路径规划、拥堵溯源和信号优化等工具，拓宽智能体的交通管理能力边界。第三，优化系统提示词设计，通过显式锁定首次出现的绝对日期来解决多工具链中的时间参照歧义。第四，探索更大规模的模型或多模型协作架构，提升系统在极端复杂交通场景下的推理鲁棒性。
